\section{Introduction}\label{sec:intro}

\subsection{Problem and research questions}
Climate change affects nutrition through two channels. The first is how much of each crop is produced. The
second is what that crop contributes to diets. A crop matters for nutrition security when three things
coincide: its growing areas are under heat or water stress, its yield falls quickly under that stress, and many
people depend on it for energy or micronutrients. This report asks:
\begin{enumerate}[label=\textbf{RQ\arabic*}]
\item Which of nine major food crops (three cereals, three pulses, three oilseeds) is most vulnerable to
observed (2001--2024) and projected (2041--2060) heat and water stress, given its nutritional importance, and
does the answer differ between countries and crop groups?
\item If a crop's yield falls, by a standard amount or by the amount implied by the projected climate change,
what happens to a country's supply of energy, protein, fat, zinc and iron, and to its nutrient adequacy?
\item What does the index indicate for one place studied in more detail: India, Telangana and the city of
Hyderabad?
\item How well supported are the crop parameters, do published future-impact studies point the same way for each
crop, what does an elevated-CO$_2$ change in nutrient concentration add to the yield channel, and how do crop losses
change the diversity of a country's food sources? (Four research extensions, Sections~\ref{sec:parevidence}--\ref{sec:divloss}.)
\end{enumerate}

\subsection{Literature context}
\paragraph{Climate and yield.} The yield literature is large and quantitative.
\begin{itemize}
\item Four independent methods agree that each degree of global warming lowers global yields of wheat, rice,
maize and soybean by about 6.0, 3.2, 7.4 and 3.1\% without CO$_2$ fertilisation or adaptation \citep{zhao2017}.
\item Year-to-year climate variation explains about a third of global yield variability \citep{ray2015}.
\item Crop models project large regional losses under warming \citep{rosenzweig2014,jagermeyr2021}.
\item Statistical studies attribute part of the recent yield trend to warming \citep{lobell2007,lobell2011}.
\item Process studies locate most of the damage in the reproductive phase, where a few hot days above a crop
threshold reduce seed set \citep{jagadish2007,deryng2014,boote2005}.
\end{itemize}
These studies measure yield; they do not ask who eats the crop.

\paragraph{Climate and nutrition.} The nutrition literature has a different emphasis.
\begin{itemize}
\item Elevated CO$_2$ lowers zinc, iron and protein in C$_3$ grains and legumes \citep{myers2014,smith2018,zhu2018,beach2019}.
\item Diet-level work asks how environmental change alters food supply and diet quality
\citep{scheelbeek2018,nelson2018,medek2017}.
\item The diversity of national food supply is itself an indicator of nutritional resilience \citep{remans2014}.
\end{itemize}

\paragraph{Vulnerability indices.} These indices usually follow the IPCC framing of exposure, sensitivity and
adaptive capacity \citep{ipcc2022ch5}, and crop-suitability studies map where crops will remain viable
\citep{ramirez2017,rippke2016}.

\paragraph{The gap.} Yield losses, nutrient supply and vulnerability scores are reported on incompatible scales.
We found no reviewed study that ranks crops by combining their climatic exposure, their sensitivity and their
nutritional importance in one index.

\subsection{Contribution}
The project contributes:
\begin{itemize}
\item a transparent index that combines the three in one number per crop and country (Section~\ref{sec:index});
\item its application to nine crops in three groups, for the present, a historical baseline and a mid-century
future, with one processing chain and one normalisation reference (Section~\ref{sec:periods});
\item an accounting step from yield loss to national nutrient supply and adequacy (Section~\ref{sec:nutmethod});
\item a local reading of the results for India and Hyderabad, using official survey data at the resolutions
they are actually published (Section~\ref{sec:india});
\item a reproducible evidence trail: raw-data checksums, a complete pipeline, generated tables, and a parameter
ledger that labels every derived value and proxy (Section~\ref{sec:repro}, Appendix~\ref{app:ledger});
\item four research extensions that test the result rather than extend its claims: an audit of how well each
parameter is supported with an evidence-screened view (Section~\ref{sec:parevidence}); a coverage assessment of
published future-impact studies for all nine crops (Section~\ref{sec:futurelit}); an elevated-CO$_2$ nutrient-quality
stress test (Section~\ref{sec:co2}); and the diversity of national food supply after crop losses
(Section~\ref{sec:divloss}).
\end{itemize}
The index is a \emph{screening} tool. It shows where stress and dependence coincide. It does not forecast yields
or diets.

\paragraph{Framing.} The structure follows the IPCC AR4 vulnerability framing (exposure $\times$ sensitivity
$\times$ adaptive capacity), with adaptive capacity replaced by nutritional \emph{importance}. In AR5/AR6
vocabulary the index is closer to a risk index, in which the stress plays the role of hazard and nutrition is
what is exposed. The substitution is deliberate, because the project asks about nutrition security, not about
farmers' coping capacity. It also means the index is not a pure climate-impact measure.

\subsection{Relation to the mid-evaluation and how to read this report}
The mid-evaluation presented the method for four crops (rice, wheat, maize, soybean). Its deck, report, tables
and figures are frozen in \file{archive/mid\_evaluation/}, and no number from them is reused as a new result.
This report completes the commitments made then.
\begin{itemize}
\item The presented \emph{literature-only} sensitivity is now used everywhere. The mid-evaluation numbers had
used an average of literature and observed-yield sensitivity; that version survives only as a robustness test.
\item The analysis covers nine crops, the future period, the nutrition consequences, and robustness,
validation and literature comparison under the final configuration.
\end{itemize}

\paragraph{Where to find things.}
\begin{itemize}
\item \textbf{Method:} Sections~\ref{sec:scope}--\ref{sec:periods}, with symbols in Table~\ref{tab:symbols}.
\item \textbf{Worked examples:} Section~\ref{sec:worked}.
\item \textbf{Results:} Sections~\ref{sec:present}--\ref{sec:india}.
\item \textbf{Research extensions} (parameter evidence, future literature, CO$_2$ quality, diversity after losses):
Sections~\ref{sec:parevidence}--\ref{sec:divloss}.
\item \textbf{Robustness, validation and literature comparison:} Section~\ref{sec:checks}.
\item \textbf{Limitations and conclusions:} Sections~\ref{sec:limits}--\ref{sec:discussion}.
\item \textbf{Reproduction and audit:} Section~\ref{sec:repro}.
\item \textbf{Appendices:} the parameter and source ledger, full result tables, the record of corrections, and
the audit evidence.
\end{itemize}

Every number from our own results is a macro or table written by the code from the saved results
(\file{docs/nine\_crop/generated/}); no result is typed by hand. Numbers quoted from published studies are typed
from the sources recorded in the evidence files and carry a citation.
